% Generated from locale/tr/content/set-theory/story/extensionality.tex; do not hand-edit.


\olfileid[tr]{sth}{story}{extensionality}
\olsection{Kaplamsallık}

Söylenmesi gereken ilk şey, kümelerin !!{element}s{}ınca bireyleştirilmesidir.
Daha kesin olarak:

\begin{axiom}[Kaplamsallık]
$A$ ve $B$ kümeleri aynı !!{element}s{}ına sahipse $A$ ile $B$ aynı kümedir.
\[
  \lforall[A][\lforall[B][(\lforall[x][(x \in A \liff x \in B)] \lif
  \eq[A][B])]]
\]
\end{axiom}

Bunu \olref[sfr][][]{part} boyunca varsaydık. Fakat yinelemeye değer.
Kaplamsallık Aksiyomu, bir kümenin !!{element}s{}ınca belirlendiği temel
düşüncesini ifade eder. (Dolayısıyla kümeler, tam olarak aynı nesnelerin pek
çok farklı kavramın kapsamına girebildiği \emph{kavramlarla} karşılaştırılabilir.)

Bu ilkeyi neden benimseyelim? Kaplamsallığın reddinin, \emph{küme teorisi}
diye adlandırılabilecek her şeyi terk etme kararı olduğunu söylemek
akla yatkındır. Küme teorisi, kaplamsal koleksiyonların teorisinden ne daha
fazla ne daha azdır.

Ne var ki \olref[sth][][]{part} içindeki asıl güçlük, \emph{hangi kümelerin
var olduğunu} söyleyen ilkeler ortaya koymaktır. Bu soruya gerçekten
``apaçık'' görünen tek yanıtın kanıtlanabilir biçimde yanlış olduğu ortaya
çıkar.

